\section{おわりに}
\label{sec_conclusion}
本稿では，大喜利投稿サイト「ボケて」の約9,250万件のボケに，構造を表すラベルを付与する方法を検討した経緯を報告した．
実例の読解から，制約の順守，発話の型，画像外の知識の3つの観点を定め，定義書を3回改訂した．
検証データ620件の付け直しにより，発話の型の6区分で F1 0.80，画像外の知識で0.85という目視ラベルの再現性を示し，これを自動判定の実質的な上限とした．
制約の順守は，カルタでは書き出しの規則で，穴埋めとお題文では同じお題の他のボケとの類似度で判定できることを示した．
発話の型と画像外の知識については，お題画像とボケを入力とする小さな VLM を，教師ラベルで学習した後に少量の目視ラベルで微調整する二段階の学習が最も有効であり，最終評価用データで6区分のマクロ F1 0.71〜0.74，3区分0.86〜0.88，画像外の知識の F1 0.68〜0.70を得た．
さらに，同じお題のボケでプレフィックスの計算を共有することで，推論を約4倍に高速化し，Qwen3-VL-4B で全件を約8日で処理できる見込みを得た．

今後は，人の判断者によるラベルで目視ラベルの妥当性を確かめるとともに，ツッコミと心の声の定義を見直す．
また，画像に写っている作品の同定と，より大きなモデルによる教師ラベルの改善により，画像外の知識の判定を改善する．
そのうえで，全件にラベルを付与し，大喜利の生成と判定の研究に用いる．
